\begin{lemma}[Off-diagonal kernel representation of the $L^1$ extension]
Fix $n \geq 1$ and $j \in \{1,\ldots,n\}$, let $\lambda$ be Lebesgue measure on
$\mathbb{R}^n$, let $\gamma_n$ be the Gaussian measure of
\noderef{GaussMeasure}, and let $K_{\gamma,j}$ be the kernel of
\noderef{GaussRieszKernel}.  Let
\[
  T : L^1(\gamma_n) \longrightarrow
      \mathcal{M}\bigl(\gamma_n;\mathbb{R}^n\bigr)
\]
be \emph{any} linear map into the $\gamma_n$-almost-everywhere equivalence
classes of measurable $\mathbb{R}^n$-valued functions such that
\begin{enumerate}
\item[(T1)] $T(\iota v) = R_\gamma v$ for every $v \in L^2(\gamma_n)$, where
$\iota : L^2(\gamma_n) \to L^1(\gamma_n)$ is the inclusion of
\noderef{ToL1} and $R_\gamma$ is the Riesz transform of
\noderef{RieszTransform}; and
\item[(T2)] $t\,\gamma_n\bigl(\{x : |Tf(x)|_{\mathbb{R}^n} > t\}\bigr)
\leq 2\,\|f\|_{L^1(\gamma_n)}$ for every $f \in L^1(\gamma_n;\mathbb{R})$ and
every real $t > 0$, the distribution function being that of the Euclidean norm
of the vector $Tf(x)$.
\end{enumerate}
Let $g \in L^1(\gamma_n)$, write $g$ also for its canonical representative, let
$x_0 \in \mathbb{R}^n$ and $\rho > 0$, and assume that
\[
  g = 0 \quad \gamma_n\text{-almost everywhere on } B(x_0,\rho) .
\]
Then:
\begin{enumerate}
\item for \emph{every} $z \in B(x_0,\rho/2)$ the function
$y \mapsto K_{\gamma,j}(z,y)\,g(y)$ is Lebesgue integrable on $\mathbb{R}^n$;
and
\item for $\gamma_n$-almost every $z$, if $z \in B(x_0,\rho/2)$ then
\[
  (Tg)_j(z) = \int_{\mathbb{R}^n}K_{\gamma,j}(z,y)\,g(y)\,dy .
\]
\end{enumerate}

\emph{The hypothetical style, and why it is not a citation of hypotheses.}
(T1) and (T2) are written out in full above; they are not imported by reference
from any other node, and the lemma is a statement about every $T$ satisfying
them.  This is the same style as \noderef{CoordinateEndpointBound}, with which
the lemma composes directly: an actual such $T$ exists by
\noderef{L1Extension}, and every consumer that has one in hand can feed it to
both nodes at once.

\emph{What this adds, and what it is not.}  The reference paper's own
interchange of the time and space integrals is licensed only ``for compactly
supported smooth inputs and $x$ off their support''.  Clause (2) upgrades that
to an arbitrary $L^1(\gamma_n)$ input, at the price of holding $z$ in a ball on
which the input vanishes and of weakening ``every $z$'' to ``almost every
$z$'' --- the latter being unavoidable, since $Tg$ is only an almost-everywhere
equivalence class.  It is emphatically \emph{not} a principal-value
identification: no limit in the truncation radius occurs, because on the range
considered the truncated integrals are eventually independent of the radius.

\emph{Clause (1) is stated and is not an afterthought.}  The integral in clause
(2) is a Bochner integral, whose value is $0$ by convention on a
non-integrable integrand; without clause (1) the identity could be read as
holding at points where the right-hand side is that conventional value.  Clause
(1) also holds at \emph{every} point of the ball, not merely almost every one,
which is what lets a consumer evaluate the right-hand side pointwise.
\end{lemma}
\begin{proof}
Throughout, $B(a,s)$ and $\overline B(a,s)$ are the open and closed Euclidean
balls, $\lambda$ is Lebesgue measure, and we write $\|h\|_1 = \int|h|\,d\gamma_n$
for a $\gamma_n$-integrable function $h$ and also for the norm of the
corresponding element of $L^1(\gamma_n)$; the two agree, because the canonical
representative of an element of $L^1(\gamma_n)$ realises its norm as that
integral.  Set
\[
  R = \|x_0\| + \rho + 1 > 0 ,
\]
so that every $z \in B(x_0,\rho/2)$ has $|z| \leq \|x_0\| + \rho/2 \leq R$, and
let $C > 0$ be the constant supplied by clause (1) of
\noderef{GaussRieszKernelFarMajorant} for the pair $(R,\rho/4)$, both of which
are positive.  For $z \in \mathbb{R}^n$ write
$\Omega_z = \{y : |z-y| > \rho/4\}$.

\emph{Step 0: the vanishing holds $\lambda$-almost everywhere as well.}
The hypothesis says that the set $N = \{y \in B(x_0,\rho) : g(y) \neq 0\}$ is
contained in a $\gamma_n$-null set.  By clause (4) of
\noderef{GaussLebesgueBallComparison}, $\gamma_n$ and $\lambda$ have the same
null sets, so $N$ is contained in a $\lambda$-null set too; that is, $g = 0$
$\lambda$-almost everywhere on $B(x_0,\rho)$.

\emph{Step 1: two separation estimates.}
Let $z \in B(x_0,\rho/2)$.
\begin{itemize}
\item[(a)] If $y \notin B(x_0,\rho)$ then
$|z-y| \geq |x_0-y| - |x_0-z| \geq \rho - \rho/2 = \rho/2 > \rho/4$, so
$y \in \Omega_z$.  Equivalently, $y \notin \Omega_z$ implies
$y \in B(x_0,\rho)$.
\item[(b)] If $|z-y| \leq \rho/4$ then
$|y-x_0| \leq |y-z| + |z-x_0| < \rho/4 + \rho/2 = 3\rho/4$, so
$y \in B(x_0,3\rho/4) \subseteq B(x_0,7\rho/8)$.
\end{itemize}

\emph{Step 2: clause (1), and the truncated form of the integral.}
We prove the following statement, which we will use both for $g$ and for the
approximants of Step 4.  \emph{Let $h : \mathbb{R}^n \to \mathbb{R}$ be
$\gamma_n$-integrable and suppose $h = 0$ $\lambda$-almost everywhere on
$B(x_0,3\rho/4)$.  Then for every $z \in B(x_0,\rho/2)$ the function
$y \mapsto K_{\gamma,j}(z,y)h(y)$ is Lebesgue integrable on $\mathbb{R}^n$ and}
\begin{equation*}
  \int_{\mathbb{R}^n}K_{\gamma,j}(z,y)h(y)\,dy
    = \int_{\Omega_z}K_{\gamma,j}(z,y)h(y)\,dy . \tag{1}
\end{equation*}
Indeed, fix such an $h$ and such a $z$.  Being $\gamma_n$-integrable, $h$
satisfies the hypothesis of clause (1)(b) of
\noderef{GaussRieszKernelFarMajorant}, and $|z| \leq R$, so that clause gives
that $y \mapsto K_{\gamma,j}(z,y)h(y)$ is Lebesgue integrable on $\Omega_z$.
By Step 1(b), every $y \notin \Omega_z$ --- that is, every $y$ with
$|z-y| \leq \rho/4$ --- lies in $B(x_0,3\rho/4)$, where $h$ vanishes
$\lambda$-almost everywhere.  Hence
$K_{\gamma,j}(z,\cdot)h$ vanishes $\lambda$-almost everywhere on the complement
of $\Omega_z$.  A function integrable on a measurable set and vanishing almost
everywhere off it is integrable on the whole space, with the same integral;
this gives both the integrability on $\mathbb{R}^n$ and $(1)$.

Applying this with $h = g$, which is $\gamma_n$-integrable and vanishes
$\lambda$-almost everywhere on $B(x_0,\rho) \supseteq B(x_0,3\rho/4)$ by
Step 0, proves clause (1) of the lemma.

\emph{Step 3: the smooth off-diagonal identification, derived inline.}
Let $\varphi : \mathbb{R}^n \to \mathbb{R}$ be smooth with compact support, let
$z \in \mathbb{R}^n$, and suppose $\varphi(y) = 0$ whenever $|z-y| \leq \sigma$,
for some $\sigma > 0$.  Write
\[
  \Phi_j^\varphi(z) = \frac{1}{\sqrt\pi}\int_0^\infty
      e^{-s}\,\frac{1}{\sqrt2}\bigl(A_s(\partial_j\varphi)\bigr)(z)\,
      \frac{ds}{\sqrt s},
  \qquad
  (A_s w)(z) = \int w\bigl(e^{-s}z + \sqrt{1-e^{-2s}}\,\zeta\bigr)d\gamma_n(\zeta) ,
\]
for the subordination integral appearing on the right-hand side of
\noderef{PVCoreSmooth} and of \noderef{RieszSemigroupPointwise}.  We claim
\begin{equation*}
  y \mapsto K_{\gamma,j}(z,y)\varphi(y) \text{ is Lebesgue integrable on }
  \mathbb{R}^n, \quad\text{and}\quad
  \Phi_j^\varphi(z) = \int_{\mathbb{R}^n}K_{\gamma,j}(z,y)\varphi(y)\,dy .
  \tag{2}
\end{equation*}
Indeed, let $0 < \varepsilon \leq \sigma$.  Clause (1) of
\noderef{PVCoreSmooth}, applied to $\varphi$, $j$, $z$ and this $\varepsilon$,
says that $y \mapsto K_{\gamma,j}(z,y)\varphi(y)$ is Lebesgue integrable on
$\{y : |z-y| > \varepsilon\}$.  Its complement is contained in
$\{y : |z-y| \leq \sigma\}$, where $\varphi$ vanishes identically, so --- by the
same elementary fact used in Step 2 --- the function is integrable on
$\mathbb{R}^n$ and
\[
  \int_{\{|z-y|>\varepsilon\}}K_{\gamma,j}(z,y)\varphi(y)\,dy
    = \int_{\mathbb{R}^n}K_{\gamma,j}(z,y)\varphi(y)\,dy
  \qquad (0 < \varepsilon \leq \sigma) .
\]
Thus the function $\varepsilon \mapsto
\int_{\{|z-y|>\varepsilon\}}K_{\gamma,j}(z,y)\varphi(y)\,dy$ is constant on
$(0,\sigma]$, so its limit as $\varepsilon \downarrow 0$ exists and equals that
constant.  Clause (2) of \noderef{PVCoreSmooth} identifies that limit as
$\Phi_j^\varphi(z)$.  This is $(2)$.

\emph{Step 4: the approximating sequence.}
First fix a cut-off.  Since $0 < 7\rho/8 < 15\rho/16$, there is a smooth bump
function $b : \mathbb{R}^n \to [0,1]$ with
$b = 1$ on $\overline B(x_0,7\rho/8)$ and
$\operatorname{supp} b \subseteq B(x_0,15\rho/16)$: this is the standard radial
bump at centre $x_0$ with inner radius $7\rho/8$ and outer radius $15\rho/16$,
which exists for any pair of radii $0 < r_{\mathrm{in}} < r_{\mathrm{out}}$.
Set $\chi = 1 - b$, so that
\[
  \chi \text{ is smooth},\quad 0 \leq \chi \leq 1,\quad
  \chi = 0 \text{ on } \overline B(x_0,\tfrac{7\rho}{8}),\quad
  \chi = 1 \text{ off } B(x_0,\tfrac{15\rho}{16}) ,
\]
the last because $b$ vanishes off its support.
Let $k \in \mathbb{N}$.  By \noderef{SmoothCompactDenseL1Gauss} applied to $g$
and $\delta = 1/(k+1) > 0$ there is a smooth, compactly supported,
$\gamma_n$-integrable $\varphi_k$ with $\|g-\varphi_k\|_1 < 1/(k+1)$; that node
states the bound for the $L^1(\gamma_n)$ seminorm read in $[0,+\infty]$, and the
bound in particular makes it finite, so $g - \varphi_k$ is $\gamma_n$-integrable
and the seminorm is the ordinary integral $\int|g-\varphi_k|\,d\gamma_n$.  Put
$\psi_k = \chi\,\varphi_k$.  Then:
\begin{itemize}
\item $\psi_k$ is smooth, being a product of smooth functions, and
$\{\psi_k \neq 0\} \subseteq \{\varphi_k \neq 0\}$, so
$\operatorname{supp}\psi_k \subseteq \operatorname{supp}\varphi_k$ is compact;
\item $\psi_k = 0$ on $\overline B(x_0,7\rho/8)$, hence in particular
$\lambda$-almost everywhere on $B(x_0,3\rho/4)$;
\item $\psi_k$ is bounded (continuous with compact support) and $\gamma_n$ is a
probability measure, hence finite (\noderef{GaussIsProb}), so $\psi_k$ is
$\gamma_n$-integrable and also square $\gamma_n$-integrable.
\end{itemize}
Moreover
\[
  \|g - \psi_k\|_1
    \leq \int|g - \chi g|\,d\gamma_n + \int|\chi|\,|g - \varphi_k|\,d\gamma_n
    \leq \int\bigl|(1-\chi)g\bigr|\,d\gamma_n + \|g-\varphi_k\|_1 ,
\]
using $0 \leq \chi \leq 1$.  The first term vanishes: $1-\chi$ is $0$ off
$B(x_0,15\rho/16)$, so $(1-\chi)g$ vanishes off $B(x_0,15\rho/16) \subseteq
B(x_0,\rho)$, and on $B(x_0,\rho)$ we have $g = 0$ $\gamma_n$-almost everywhere
by hypothesis; so $(1-\chi)g = 0$ $\gamma_n$-almost everywhere.  Therefore
\begin{equation*}
  \|g-\psi_k\|_1 < \frac{1}{k+1} . \tag{3}
\end{equation*}
Write $u_k \in L^2(\gamma_n)$ for the class of $\psi_k$ and
$\iota\psi_k = \iota u_k \in L^1(\gamma_n)$ for its image under the inclusion of
\noderef{ToL1}; by the definition of that inclusion the canonical
representative of $\iota\psi_k$ equals $\psi_k$ $\gamma_n$-almost everywhere,
so $\|g - \iota\psi_k\|_{L^1(\gamma_n)} = \|g-\psi_k\|_1 < 1/(k+1)$.

\emph{Step 5: the right-hand side converges, uniformly on the ball.}
Let $z \in B(x_0,\rho/2)$.  By Step 4 the function $\psi_k$ satisfies the
hypotheses of Step 2, so $(1)$ holds for it; and $(1)$ holds for $g$ by Step 2.
Both $K_{\gamma,j}(z,\cdot)\psi_k$ and $K_{\gamma,j}(z,\cdot)g$ are integrable
on $\Omega_z$, so by linearity of the integral there and by clause (1)(b) of
\noderef{GaussRieszKernelFarMajorant} applied to the $\gamma_n$-integrable
function $\psi_k - g$ at the point $z$ (recall $|z| \leq R$),
\begin{equation*}
  \Bigl|\int_{\mathbb{R}^n}K_{\gamma,j}(z,y)\psi_k(y)\,dy
        - \int_{\mathbb{R}^n}K_{\gamma,j}(z,y)g(y)\,dy\Bigr|
  = \Bigl|\int_{\Omega_z}K_{\gamma,j}(z,y)\bigl(\psi_k-g\bigr)(y)\,dy\Bigr|
  \leq C\,\|\psi_k - g\|_1 < \frac{C}{k+1} , \tag{4}
\end{equation*}
the bound being independent of $z \in B(x_0,\rho/2)$.

\emph{Step 6: the left-hand side converges in measure.}
Let $t > 0$ and $k \in \mathbb{N}$.  Applying
\noderef{CoordinateEndpointBound} --- to this $T$, with the two properties it
requires, which are (T1) and (T2) and hold by assumption --- to the element
$g - \iota\psi_k$ of $L^1(\gamma_n)$, to the coordinate $j$ and to $t$,
\[
  t\,\gamma_n\bigl(\{z : |\,(T(g - \iota\psi_k))_j(z)\,| > t\}\bigr)
    \leq 2\,\|g-\iota\psi_k\|_{L^1(\gamma_n)} < \frac{2}{k+1} .
\]
Since $T$ is linear, $T(g-\iota\psi_k) = Tg - T(\iota\psi_k)$ as
almost-everywhere classes, and the canonical representative of a difference of
such classes is $\gamma_n$-almost everywhere the difference of the
representatives.  The two sets
$\{z : |(T(g-\iota\psi_k))_j(z)| > t\}$ and
$\{z : |(Tg)_j(z) - (T(\iota\psi_k))_j(z)| > t\}$ therefore differ inside a
$\gamma_n$-null set and have the same $\gamma_n$-measure, so
\[
  t\,\gamma_n\bigl(\{z : |\,(Tg)_j(z) - (T(\iota\psi_k))_j(z)\,| > t\}\bigr)
    < \frac{2}{k+1} .
\]
Hence for each fixed $t>0$ the measures
$\gamma_n(\{|(Tg)_j - (T(\iota\psi_k))_j| > t\})$ tend to $0$ as
$k \to \infty$; that is, $(T(\iota\psi_k))_j \to (Tg)_j$ in
$\gamma_n$-measure.  Since $\gamma_n$ is a finite measure
(\noderef{GaussIsProb}), a sequence converging in measure has a subsequence
converging almost everywhere; fix $k_1 < k_2 < \cdots$ with
\begin{equation*}
  (T(\iota\psi_{k_m}))_j(z) \longrightarrow (Tg)_j(z)
  \qquad\text{for } \gamma_n\text{-almost every } z . \tag{5}
\end{equation*}

\emph{Step 7: identification of each approximant, and the conclusion.}
Fix $m$.  The function $\psi_{k_m}$ is smooth with compact support, and for
every $z \in B(x_0,\rho/2)$ it vanishes on $\{y : |z-y| \leq \rho/4\}$, because
by Step 1(b) such a $y$ lies in $B(x_0,7\rho/8)$ where $\psi_{k_m}$ is zero.  So
Step 3 applies with $\sigma = \rho/4$ and gives
\begin{equation*}
  \Phi_j^{\psi_{k_m}}(z)
    = \int_{\mathbb{R}^n}K_{\gamma,j}(z,y)\psi_{k_m}(y)\,dy
  \qquad\text{for every } z \in B(x_0,\rho/2) . \tag{6}
\end{equation*}
On the other hand, clause (3) of \noderef{RieszSemigroupPointwise}, applied to
the smooth compactly supported $\psi_{k_m}$ (which is square
$\gamma_n$-integrable by Step 4) and to $j$, states that
$\Phi_j^{\psi_{k_m}}$ agrees $\gamma_n$-almost everywhere with the $j$th
coordinate of $R_\gamma u_{k_m}$, $u_{k_m}$ being the $L^2(\gamma_n)$ class of
$\psi_{k_m}$.  By (T1), $T(\iota u_{k_m}) = R_\gamma u_{k_m}$ as
almost-everywhere classes, so their representatives agree $\gamma_n$-almost
everywhere.  Combining,
\begin{equation*}
  (T(\iota\psi_{k_m}))_j(z) = \Phi_j^{\psi_{k_m}}(z)
  \qquad\text{for } \gamma_n\text{-almost every } z . \tag{7}
\end{equation*}

Now intersect countably many sets of full $\gamma_n$-measure: the one from
$(5)$, and for each $m$ the one from $(7)$.  A countable intersection of
conull sets is conull.  Let $z$ belong to it and suppose
$z \in B(x_0,\rho/2)$.  By $(7)$ and $(6)$,
$(T(\iota\psi_{k_m}))_j(z) = \int_{\mathbb{R}^n}K_{\gamma,j}(z,y)\psi_{k_m}(y)\,dy$
for every $m$.  Letting $m \to \infty$, the left-hand side tends to $(Tg)_j(z)$
by $(5)$, and the right-hand side tends to
$\int_{\mathbb{R}^n}K_{\gamma,j}(z,y)g(y)\,dy$ by $(4)$, since
$C/(k_m+1) \to 0$.  A convergent sequence of reals has a unique limit, so
\[
  (Tg)_j(z) = \int_{\mathbb{R}^n}K_{\gamma,j}(z,y)\,g(y)\,dy .
\]
As this holds for every $z$ in a set of full $\gamma_n$-measure, clause (2)
follows.
\end{proof}
